\documentclass[10pt,a4paper]{amsart}
\usepackage[margin=19mm]{geometry}
\usepackage{amsmath,amssymb,mathtools}
\usepackage[scr=rsfs,cal=euler]{mathalfa}
\usepackage{xfrac}
\usepackage[unicode,hidelinks,bookmarks=false]{hyperref}
\newcommand{\ie}{i.e.\ }
\newcommand{\cf}{cf.\ }
\newcommand{\resp}{respectively\ }
\newcommand{\Subsection}{\subsection}
\providecommand{\nobreakdash}{\nobreak\hbox{-}}
\setlength{\emergencystretch}{2em}
\multlinegap0pt
\usepackage{bm}
\usepackage{bbm}
\usepackage{dsfont}
\usepackage{xspace}
\usepackage{cancel}
\usepackage{mathtools}
\usepackage[shortlabels]{enumitem}
\usepackage{amsthm}
\makeatletter
\@ifundefined{theorem}{\newtheorem{theorem}{Theorem}}{}
\@ifundefined{claim}{\newtheorem{claim}[theorem]{Claim}}{}
\@ifundefined{lemma}{\newtheorem{lemma}[theorem]{Lemma}}{}
\@ifundefined{proposition}{\newtheorem{proposition}[theorem]{Proposition}}{}
\makeatother
\def\Z{{\mathbb{Z}}}
\def\n{{\mathrm{N}}}
\newcommand{\D}{\Delta}
\newcommand{\cil}{\lceil\frac{r}{2}\rceil}
\newcommand{\f}{\lfloor\frac{r}{2}\rfloor}
\newcommand{\lk}{\mbox{\upshape lk}\,}
\title[On the Vietoris-Rips Complexes of Integer Lattices]{On the Vietoris-Rips Complexes of Integer Lattices}
\author{Raju Kumar Gupta, Sourav Sarkar, Samir Shukla}
\date{Source: arXiv:2511.04238v1 (CC BY 4.0)}
\begin{document}
\maketitle

\noindent\emph{Source and licence.} On the Vietoris-Rips Complexes of Integer Lattices. Version arXiv:2511.04238v1 (CC BY 4.0), distributed under the Creative Commons Attribution 4.0 International licence, as recorded in the arXiv licence metadata for this exact version and shown on the abstract page https://arxiv.org/abs/2511.04238v1. Full rights notice: licenses/arxiv-2511.04238-CC-BY-4.0.txt.\par\smallskip

\noindent\textit{Editorial scope.} Counted unit: the Lemma whose source label is \texttt{sumoffour2} in the article (source0.tex lines 805--812, source label \texttt{sumoffour2}), delivered together with the article's own complete original proof of it (source0.tex lines 814--961, one \texttt{proof} environment of 148 lines). No mathematics is written, repaired or re-derived: statement and proof are the article's own text line for line. Required context retained uncounted: flag (lines 276--278); threevertex (lines 637--645); sumofthree (lines 663--665); fourvertex (lines 737--745); sumoffour (lines 764--766). Every definition, display and label of those lines is retained unchanged. Omitted from this package and not redistributed: 1--275, 279--636, 646--662, 666--736, 746--763, 767--804, 813--813, 962--2150. Nothing inside a retained excerpt is omitted or altered: every retained body is delivered exactly as the article's lines are written, so no omission receipt of any kind accompanies this package. Citations: the retained excerpts carry no citation command at all, so no bibliography entry of the article is transcribed and no reference is redistributed. Reference adaptation: none was needed and none was made; every reference inside the delivered unit resolves inside it, so no unresolved reference or citation is produced. Printed slips kept literal: no printed word, symbol, hypothesis or number of the counted statement or of its proof was altered. Layout and class: the article's own class is \texttt{amsart} with packages including amsmath, amssymb, amsthm, mathtools, amsfonts, bbm, inputenc, xy, graphicx, epsfig, multirow, enumerate, xcolor, color; the wrapper is the lane's pinned \texttt{amsart} editorial wrapper at 10pt on A4, which restates the theorem-like environments the retained text needs and adds the article's own macro definitions that the retained text needs (\texttt{\textbackslash D}, \texttt{\textbackslash Z}, \texttt{\textbackslash cil}, \texttt{\textbackslash f}, \texttt{\textbackslash lk}, \texttt{\textbackslash n}), with extra wrapper packages (bm, bbm, dsfont, xspace, cancel, mathtools, enumitem). The delivered numbering is the unit's own. Assets: no retained span contains a figure environment, a graphics-inclusion command, a PDF-page inclusion, a table or a TikZ picture; no image file is copied into this package, the package declares no asset dependency and no image file is redistributed with it. Licence: version arXiv:2511.04238v1 of this article is distributed under the Creative Commons Attribution 4.0 International licence, as recorded in the arXiv licence metadata for this exact version and shown on the abstract page https://arxiv.org/abs/2511.04238v1. A full rights notice is delivered at licenses/arxiv-2511.04238-CC-BY-4.0.txt. Characters: every retained excerpt is pure ASCII, so the delivered engine, XeLaTeX with its default Latin Modern text font, reports no missing character.
\begin{proposition}\label{flag}
	Let $K$ be a simplicial complex and $a,b\in V(K)$ such that $a \neq b$. Let $\n[a, K]\subseteq \n[b, K]$. If $K$ is a flag complex, then $K \simeq K \setminus a$.
\end{proposition}

\begin{lemma}\label{threevertex}
	Let $r\geq 2$, and let  $\D$ be  a subcomplex of $\Gamma_n^{\alpha, r}$ that contains the vertex $x$. Then for any two vertices $y$ and $z$ in $\D$, where $z\in V(\lk(x,\D))$ and  $|y_l|=|x_l|-1$ for $l\in\{i,j,k\}$ and $y_l=x_l$ for $l\in [n]\setminus\{i,j,k\}$, we have the following: 
	\begin{itemize}
		\item[(i)]  If $|z_l|<|x_l|$ for at least two choices of $l\in\{i,j,k\}$,  then $d(z,y)\leq  r-1$.
		\item[(ii)]  Let $|z_k|<|x_k|$. If $|z_i|\geq |x_i|$ and $sgn(z_ix_i)=-1$,  or $|z_j|\geq |x_j|$ and $sgn(z_jx_j)=-1$, then $d(z,y)\leq  r$.
		\item[(iii)] If $|z_l|\geq |x_l|$ for all $l\in\{i,j,k\}$ and $sgn(z_sx_s)=-1$ for at least two choices of $s\in\{i,j,k\}$, then $d(z,y)\leq r$.
	\end{itemize}
	
\end{lemma}

\begin{lemma}\label{sumofthree}
	Let $r\geq n\geq 4$ . Then $\Gamma_n^{\alpha, r}$ is homotopy equivalent to the induced subcomplex $\Delta'$ of $\Gamma_n^{\alpha, r}$, where every $x \in V(\Delta')$ satisfies the following: (i) $|x_i|< \f, |x_j|+|x_k|\leq \cil$ for all $i, j, k \in [n], j \neq k$, and (ii) $|x_i| + |x_j| + |x_k| < r-1$ for all distinct $i, j, k \in [n]$. 
\end{lemma}

\begin{lemma}\label{fourvertex}
	Let $r\geq 4$, and let  $\D$ be  a subcomplex of $\Gamma_n^{\alpha, r}$ that contains the vertex $x$. Then for any two $y, z \in V(\D)$, where $z\in V(\lk(x,\D))$ and  $|y_t|=|x_t|-1$ for $t\in\{i,j,k,l\}$ and $y_s=x_s$ for $s\in [n]\setminus\{i,j,k,l\}$, we have the following: 
	\begin{enumerate}[label=(\roman*)]
		\item If $|z_s|<|x_s|$ for at least two choices of $s\in\{i,j,k,l\}$,  then $d(z,y)\leq r$. 
		\item Let $|z_s|<|x_s|$ for some $s \in \{i, j, k, l\}$. If there exists $t \in \{i, j, k, l\} \setminus \{s\}$ such that $|z_t|\geq |x_t|$ and $sgn(z_tx_t)=-1$, then $d(z,y)\leq  r$. 
		\item If $|z_t|\geq |x_t|$ for all $t\in\{i,j,k,l\}$, and  $sgn(z_sx_s)=-1$  for at least two indices $s\in\{i,j,k,l\}$, then $d(z,y)\leq r$.
	\end{enumerate}
	
\end{lemma}

\begin{lemma}\label{sumoffour}
	Let $r\geq n\geq 4$ . Then $\Gamma_n^{\alpha, r}$ is homotopy equivalent to the induced subcomplex $\Delta'$ of $\Gamma_n^{\alpha, r}$, where every $x \in V(\Delta')$ satisfies the following: (i) $|x_i|< \f, |x_j|+|x_k|\leq \cil$ for all $i, j, k \in [n], j \neq k$ (ii) $|x_i| + |x_j| + |x_k| < r-1$ for all $\{i, j, k\}\subseteq [n]$, and (iii) $|x_i| + |x_j| + |x_k| + |x_l|\leq r-1$ for all  $\{i, j, k, l\} \subseteq [n]$.
\end{lemma}

\begin{lemma}\label{sumoffour2}
	Let $r \geq n \geq 5$ and  $r\geq 10$. Then $\Gamma_n^{\alpha, r}$ is homotopy equivalent to the induced subcomplex $\Delta$, where $x \in V(\Delta)$ satisfies the following:
	(i) $|x_i|< \f, |x_j|+|x_k|\leq \cil$ for all $i, j, k \in [n], j \neq k$ and (ii) there exists no $\{i, j, k , l\}\subseteq [n]$ such that $|x_i|+ |x_j|+ |x_k| + |x_l| \geq  r-1$ and $x_s=0$ for all $s\notin\{i,j,k,l\}$.
	
	
	
	
\end{lemma} 

\begin{proof}
	From Lemma \ref{sumoffour},  $\Gamma_n^{\alpha, r}$ is homotopy equivalent to the induced subcomplex $\Delta_1$ of $\Gamma_n^{\alpha, r}$, where $x \in V(\Delta_1)$ satisfies the three conditions $(i), (ii)$ and $(iii)$, as given in Lemma \ref{sumoffour}. 
	
	Let $A = \{x \in V(\Delta_1) : \exists \ \{i, j, k, l\} \subseteq [n] \ \text{such that} \  |x_i|+ |x_j|+ |x_k| + |x_l| = r-1 \ \text{and }  x_s = 0  \text{ for all }  s \notin\{i,j,k,l\} \}$. It is sufficient to show that $\D_1\simeq \mathrm{Ind}_{\D_1}(V(\D_1)\setminus \{A\})$. 
	
 Let $B=\{z\in V(\D_1): \sum_{s\in\{i,j,k,l\}}|z_s|=r-1, \ \text{and} \ i<j<k<l, z_l\geq2\ \text{for some} \ \{i,j,k,l\}\subseteq[n] 
	\ \text{and} \  z_s=0 \ \text{for}\ s\notin\{i,j,k,l\} \}$. 
	
	Let $C=\{z\in V(\D_1)\setminus B: \sum_{s\in\{i,j,k,l\}}|z_s|=r-1 \ \text{and}\  i<j<k<l, z_l=1\ \text{for some} \ \{i,j,k,l\}\subseteq[n] 
	\ \text{and} \  z_s=0 \ \text{for}\ s\notin\{i,j,k,l\} \}$.
	Observe that $A=B\cup C$.
	
	We shall remove all the elements of $A$ from $\D_1$ without changing the homotopy type of $\D_1$ in two steps: in Step I, we shall remove the elements of $B$, and then in Step II, we shall remove the elements of $C$.
	
    \smallskip
	\noindent 
	\textbf{Step I:}
	Let $x\in B$. Then, $|x_i|+|x_j|+|x_k|+|x_l|=r-1$, where $\{i,j,k,l\}\subseteq [n]$, $x_l\geq 2$, $i<j<k<l$, and $x_s=0$ for all $s\notin \{i,j,k,l\}$. If $|x_t|=0$ for some $t\in \{i,j,k,l\}$, then $\sum_{s\in\{i,j,k,l\}\setminus\{t\}}|x_s|=r-1$, which contradicts Lemma \ref{sumofthree} $(ii)$. Therefore, $|x_t|\geq 1$ for all $t\in \{i,j,k,l\}$.  
	
	Consider the vertex $\lambda^{[x;ijkl]}\in V(\Z^n)$. Clearly,  $d(x,\lambda^{[x;ijkl]})=4$, and $|x_s-\lambda_s|=1$ for all $s\in\{i,j,k,l\}$. Since $i<j<k<l$ and $\lambda^{[x;ijkl]}_l>0$, and $\lambda^{[x;ijkl]}_t=0$ for $t>l$, it follows that $\lambda^{[x;ijkl]}\succ {\bf{0}}$. On the other hand $|\lambda^{[x;ijkl]}_s|<|x_s|$ for all $s\in \{i,j,k,l\}$ implies that $\lambda^{[x;ijkl]}\in \D_1$. 
	
	\begin{claim}\label{claim:fourvertex}
		$\n[x, \D_1] \subseteq \n[\lambda^{[x;ijkl]}, \D_1]$.
	\end{claim}
	\noindent {\em Proof of Claim \ref{claim:fourvertex}:}
	Since $r \geq 5$, $x\in\n[\lambda^{[x;ijkl]}, \D_1]$. Let $y\in\n[x,\D_1]$ with $x\neq y$. Suppose $y\notin \n[\lambda^{[x;ijkl]},\D_1]$.
	Then $d(y,\lambda^{[x;ijkl]})>r$. From Lemma \ref{fourvertex} $(i)$, $|y_s|<|x_s|$ for at most one value of $s$ in $\{i,j,k,l\}$, otherwise $d(y,\lambda^{[x;ijkl]})\leq r$, which is a contadiction. We have the following cases:
	
	\smallskip 
	
	\noindent\textbf{Case 1.1:} There exists $p \in \{i, j, k, l\}$ such that $|y_p| < |x_p|$. 
	
	
	Then  $|y_s|\geq |x_s|$  for all $s\in \{ i, j,k,l\} \setminus \{p\}$. If  $sgn(y_sx_s)=-1$ for some $s \in \{i, j, k, l\} \setminus \{p\}$, then from Lemma \ref{fourvertex} $(ii)$, we get $d(y,\lambda^{[x;ijkl]})\leq r$,  which is a contradiction. Hence $sgn(y_sx_s)=1$ for all $s \in \{i, j, k, l\} \setminus \{p\}$. Let $|y_s|=|x_s|+a_s$, where $a_s\geq 0$ for $s \in \{i, j, k, l\} \setminus \{p\}$.  Since $\sum_{s}|y_s|\leq r$, we see that $\sum_{s\notin \{i, j,k,l\} \setminus \{p\}}|y_s|\leq r-(\sum_{s \in \{i, j, k, l \} \setminus \{p\}}|x_s|+a_s)$. Therefore,
	\begin{align*}
		d(y,\lambda^{[x;ijkl]})&=|y_p-\lambda^{[x;ijkl]}_p|+ \sum_{s \in \{i, j, k, l\}\setminus\{p\}}|y_s-\lambda^{[x;ijkl]}_s|+\sum_{s\notin\{i,j,k,l\}}|y_s-\lambda^{[x;ijkl]}_s| \\
		&\leq  |y_p-x_p|-1+ \sum_{s \in \{i, j, k, l\}\setminus\{p\}}(|y_s-x_s|+1 )+\sum_{s\notin\{i,j,k,l\}}|y_s| \\
		&\leq|x_p|+2+\sum_{s \in \{i, j, k, l\}\setminus \{p\}} a_s+\sum_{s\notin\{i, j,k,l\} \setminus \{p\}}|y_s| \\
		&\leq |x_p|+2 + \sum_{s \in \{i, j, k, l\}\setminus \{p\}} a_s+r- \sum_{s \in \{i, j, k, l\}\setminus \{p\}} (|x_s|+a_s) \\
		&= r+2+|x_p|-\sum_{s \in \{i, j, k, l\}\setminus \{p\}} |x_s|.
	\end{align*}
	
	
	If $|x_p|-\sum_{s \in \{i, j, k, l\}\setminus \{p\}} |x_s|\leq -2$, then $d(y,\lambda^{[x;ijkl]})\leq r$, thereby implying that $y \in \n[\lambda^{[x;ijkl]}, \D_1]$. So, assume that $|x_p|-\sum_{s \in \{i, j, k, l\}\setminus \{p\}} |x_s|\geq -1$. On the other hand if $|x_p|-\sum_{s \in \{i, j, k, l\}\setminus \{p\}} |x_s|\geq 1$, then using the fact that $|x_p|+\sum_{s \in \{i, j, k, l\}\setminus \{p\}} |x_s| =r-1$, we get that $|x_p|\geq \f$, which is not possible. Thus, we have only two possibilities, either $|x_p|=\sum_{s \in \{i, j, k, l\}\setminus \{p\}} |x_s|$ or $|x_p|+1=\sum_{s \in \{i, j, k, l\}\setminus \{p\}} |x_s|$.
	
	Now, if $|x_p|=\sum_{s \in \{i, j, k, l\}\setminus \{p\}} |x_s|$, then using equation $|x_p|+\sum_{s \in \{i, j, k, l\}\setminus \{p\}} |x_s| =r-1$, we get that $|x_p|=\frac{r-1}{2}$. This implies that $r$ must be odd and $|x_p|=\f$, which is not possible. 
	Therefore, we have $|x_p|+1=\sum_{s \in \{i, j, k, l\}\setminus \{p\}}$. In this case, $|x_p|=\frac{r-2}{2}$, which implies that $r$ is even. Since $r\geq 10$, there exists a $t\in \{i, j, k, l\}\setminus \{p\}$ such that $|x_t|\geq 2$. This implies that $|x_p|+|x_t|>\frac{r}{2}$. This is a contradiction.
	
	Thus we conclude that $d(y, \lambda^{[x;ijkl]}) \leq r$, and therefore $y \in \n[\lambda^{[x;ijkl]}, \D_1]$.
	

    \smallskip
	\noindent 
	\textbf{Case 1.2} $|y_s|\geq |x_s|$ for all $s\in\{i,j,k,l\}$. 
	
	
	From Lemma \ref{fourvertex} $(iii)$, $sgn(y_sx_s)=-1$ is possible for at most one value of $s$ in $\{i,j,k,l\}$. 
	If $|y_s|\geq |x_s|$, and $sgn(y_sx_s)=1$ for all $s\in \{i,j,k, l\}$, then from $|y_i|+|y_j|+|y_k|+|y_l|\leq r-1$, it follows that $x_s=y_s$ for all $s\in \{i,j,k,l\}$. Consequently, $\sum_{s\notin\{i,j,k,l\}}|y_s|\leq 1$ implies that $d(y,\lambda^{[x;ijkl]})\leq 5<r$, which is a contradiction.  
	
	Thus, there exists a $q\in\{i,j,k,l\}$ such that $sgn(x_qy_q)=-1$. Then $sgn(x_sy_s)=1$ for $s\in\{i,j,k,l\}\setminus\{q\}$. Since $|y_i|+|y_j|+|y_k|+|y_l|\leq r-1$, it follows that $y_q=-x_q$ and $y_s=x_s$ for $s\in\{i,j,k,l\}\setminus\{q\}$. Then, we have $\sum_{s\notin\{i,j,k,l\}}|y_s|\leq 1$. Therefore, 
	\begin{align}
		d(y,\lambda^{[x;ijkl]})&=|y_q-x_q|-1+\sum_{s\in\{i,j,k,l\}\setminus\{q\}}(|y_s-x_s|+1)+\sum_{s\notin\{i,j,k,l\}}|y_s| \notag
		\\
		&\leq |y_q|+|x_q|+3. \label{eq2}
	\end{align}
	
	If $r$ is odd then $|x_q|=|y_q|\leq \f-1$. Thus, from Equation (\ref{eq2}),
	$d(y,\lambda^{[x;ijkl]})\leq \f-1+\f-1+3=r$, which is a contradiction.  Therefore, $r$ must be even. 
	Further, if $|x_q|\leq \frac{r}{2}-2$, then  $d(y,\lambda^{[x;ijkl]})\leq \frac{r}{2}-1+\frac{r}{2}-2+3=r$, which is a contradiction. From the bound  $|x_q|<\frac{r}{2}$, we must have $|x_q|=\frac{r}{2}-1$. Since $r\geq 10$, there exists a $t\in \{i, j, k, l\}\setminus \{q\}$ such that $|x_t|\geq 2$. This implies that $|x_q|+|x_t|>\frac{r}{2}$, which is a contradiction.
	
	Thus we conclude that $d(y, \lambda^{[x;ijkl]}) \leq r$, and therefore $y \in \n[\lambda^{[x;ijkl]}, \D_1]$.  Hence, $\n[x,\D_1]\subseteq \n[\lambda^{[x;ijkl]},\D_1]$. This proves Claim \ref{claim:fourvertex}.  
	
	
	From Proposition \ref{flag}, $\D_1\simeq \mathrm{Ind}_{\D_1}(V(\D_1)\setminus \{x\})$. Now, applying the same arguments for each $b\in B$, we find that $\D_1\simeq \mathrm{Ind}_{\D_1}(V(\D_1)\setminus B)$. 
	Let $\D_2=\mathrm{Ind}_{\D_1}(V(\D_1)\setminus B)$. Then $\D_1\simeq \D_2$.
	
	
	
    \smallskip
	\noindent \textbf{Step II:}
	Let $x\in C\subseteq \D_2$ be such that $|x_i|+|x_j|+|x_k|+|x_l|=r-1$  where $\{i,j,k,l\}\subseteq [n]$, $i<j<k<l$, $x_l=1$, and $x_s=0$ for all $s\notin \{i,j,k,l\}$. If $|x_t|=0$ for some $t\in \{i,j,k,l\}$, then $\sum_{s\in\{i,j,k,l\}\setminus\{t\}}|x_s|=r-1$, which contradicts Lemma \ref{sumoffour} $(ii)$. Therefore, $|x_t|\geq 1$ for all $t\in \{i,j,k,l\}$.
	
	Consider $\lambda^{[x;ijk]}\in V(\Z^n)$. Clearly, $d(x,\lambda^{[x;ijk]})=3$, and $|x_s-\lambda^{[x;ijk]}_s|=1$ for $s\in\{i,j,k\}$. Since  $x_l=1$, and $x_s=0$ for $s>l$, it follows that $\lambda^{[x;ijk]}_l=1$ and $\lambda^{[x;ijk]}_s=0$ for $s>l$. Thus, $\lambda\succ{\bf{0}}$. Moreover, since $|\lambda^{[x;ijk]}_s|<|x_s|$ for all $s\in\{i,j,k\}$, and $\lambda^{[x;ijk]}_s=p_s$ for $s\notin\{i,j,k\}$, we have $\lambda^{[x;ijk]}\in\D_2$.
    
	\vspace{-0.5 cm}
    
	\begin{claim}\label{claim:fourvertex2}
		$\n[x, \D_2] \subseteq \n[\lambda^{[x;ijk]}, \D_2]$.
	\end{claim}
    
    \vspace{-0.5 cm}
    
	\noindent {\em Proof of Claim \ref{claim:fourvertex2}:}
	Since $r > 3$, $x\in\n[\lambda^{[x;ijk]}, \D_2]$. Let $y\in\n[x,\D_2]$ with $x\neq y$. Suppose $y\notin \n[\lambda^{[x;ijk]},\D_2]$. Then $d(y,\lambda^{[x;ijk]})>r$. From Lemma \ref{threevertex} $(i)$, $|y_s|<|x_s|$ for at most one value of $s$ in  $\{i,j,k\}$, otherwise $d(y,\lambda^{[x;ijk]})\leq r$, which is a contadiction. We have the following cases:
	
	
    \smallskip
	\noindent \textbf{Case 2.1:} There exists $p \in \{i, j, k\}$ such that $|y_p| < |x_p|$. 
	
	
	Then  $|y_s|\geq |x_s|$  for all $s\in \{i ,j,k\}\setminus \{p\}$. If  $sgn(y_sx_s)=-1$ for some $s \in \{i, j, k\} \setminus \{p\}$, then from Lemma \ref{threevertex} (ii), we get $d(y,\lambda^{[x;ijk]})\leq r$,  which is a contradiction. Hence $sgn(y_sx_s)=1$ for all $s \in \{i, j, k\} \setminus \{p\}$. Let $|y_s|=|x_s|+b_s$, where $b_s\geq 0$ for $s \in \{i, j, k\} \setminus \{p\}$.  Since $\sum_{s}|y_s|\leq r$, we see that $\sum_{s\notin \{i, j,k\} \setminus \{p\}}|y_s|\leq r-(\sum_{s \in \{i, j, k \} \setminus \{p\}}|x_s|+b_s)$.
	
	\begin{align*}
    \begin{split}
		d(y,\lambda^{[x;ijk]})&=|y_p-\lambda^{[x;ijk]}_p|+ \sum_{s \in \{i, j, k\}\setminus\{p\}}|y_s-\lambda^{[x;ijk]}_s|+\sum_{s\notin\{i,j,k\}}|y_s-\lambda^{[x;ijk]}_s| \\
		&\leq  |y_p-x_p|-1+ \sum_{s \in \{i, j, k\}\setminus\{p\}}(|y_s-x_s|+1 )+\sum_{s\notin\{i,j,k\}}|y_s| +1 \\
		&\leq  |x_p|+2+ \sum_{s \in \{i, j, k\}\setminus \{p\}} b_s+\sum_{s\notin\{i,j,k\}}|y_s| +|y_p| \\
		&\leq|x_p|+2+\sum_{s \in \{i, j, k\}\setminus \{p\}} b_s+\sum_{s\notin\{i, j,k\} \setminus \{p\}}|y_s| \\
		&\leq |x_p|+2 + \sum_{s \in \{i, j, k\}\setminus \{p\}} b_s+r- \sum_{s \in \{i, j, k\}\setminus \{p\}} (|x_s|+b_s) \\
		&= r+2+|x_p|-\sum_{s \in \{i, j, k\}\setminus \{p\}} |x_s|.
         	\end{split}
	\end{align*}
	
	If $|x_p|-\sum_{s \in \{i, j, k\}\setminus \{p\}} |x_s|\leq -2$, then $d(y,\lambda^{[x;ijk]})\leq r$, which is a contradiction. If $|x_p|-\sum_{s \in \{i, j, k\}\setminus \{p\}} |x_s|\geq 1$, then using the fact that $|x_p|+\sum_{s \in \{i, j, k\}\setminus \{p\}} |x_s|=r-2$, we find that $|x_p|\geq \f$, which is not possible. Thus, we have only two possibilities: either $|x_p|=\sum_{s \in \{i, j, k\}\setminus \{p\}} |x_s|$ or $|x_p|+1=\sum_{s \in \{i, j, k\}\setminus \{p\}} |x_s|$. If $|x_p|=\sum_{s \in \{i, j, k\}\setminus \{p\}} |x_s|$, then from the equation, $|x_p|+\sum_{s \in \{i, j, k\}\setminus \{p\}} |x_s|=r-2$, we get $|x_p|=\frac{r-2}{2}$. This implies that $r$ is even. Since $r\geq 10$, there exists $t\in \{i,j,k\}\setminus \{p\}$ such that $|x_t|\geq 2$. Then $|x_p|+|x_t|>\frac{r}{2}$, which is a contradiction. 
	
	Now, we assume that $|x_p|+1=\sum_{s \in \{i, j, k\}\setminus \{p\}} |x_s|$. Then $|x_p|=\frac{r-3}{2}$. This implies that $r$ is odd and $|x_p|=\f-1$. Since $|x_l|=1$, we get $\sum_{s \in \{i, j, k\}\setminus \{p\}} |x_s|=\f$. If $|x_t|\geq 3$ for some $t\in \{i, j, k\}\setminus \{p\}$, then $|x_p|+|x_t|>\cil$, which is a contradiction. Therefore,  $\sum_{s \in \{i, j, k\}\setminus \{p\}} |x_s|\leq 4$. This implies that $\f\leq 4$ and thus, $r\leq 9$, which is a contradiction. Thus we conclude that $d(y, \lambda^{[x;ijk]}) \leq r$, and therefore $y \in \n[\lambda^{[x;ijk]}, \D_2]$.
	
	
	
	
    \smallskip
	\noindent \textbf{Case 2.2:} $|y_s|\geq |x_s|$ for all $s\in\{i,j,k\}$. 
	
	
	From Lemma \ref{threevertex} $(iii)$, $sgn(y_sx_s)=-1$ is possible for at most one value of $s$ in $\{i,j,k\}$. 
	If $|y_s|\geq |x_s|$, and $sgn(y_sx_s)=1$ for all $s\in \{i,j,k\}$, then since from Lemma \ref{sumoffour} $(iii)$,   $|y_i|+|y_j|+|y_k|\leq r-2$, it follows that $x_s=y_s$ for all $s\in \{i,j,k\}$. Consequently, $\sum_{s\notin\{i,j,k\}}|y_s|\leq 2$ implying that $d(y,\lambda^{[x;ijk]})\leq 6$, which contradicts our assumption that $y\notin \n[\lambda^{[x;ijk]},\D_2]$.  
	
	Thus, there exists a $q\in\{i,j,k\}$ such that $sgn(x_qy_q)=-1$. Then $sgn(x_sy_s)=1$ for $s\in\{i,j,k\}\setminus\{q\}$. Since $|y_i|+|y_j|+|y_k|\leq r-2$, it follows that $y_q=-x_q$, $y_s=x_s$, for $s\in\{i,j,k\}\setminus\{q\}$. Then, we have $\sum_{s\notin\{i,j,k\}}|y_s|\leq 2$. Therefore, 
	\begin{align}
		d(y,\lambda^{[x;ijk]})&=|y_q-x_q|-1+\sum_{s\in\{i,j,k\}\setminus\{q\}}(|y_s-x_s|+1)+\sum_{s\notin\{i,j,k\}}|y_s-x_s| \notag
		\\
		&\leq |y_q|+|x_q|+1+\sum_{s\notin\{i,j,k\}}|y_s|+\sum_{s\notin\{i,j,k\}}|x_s| \notag
		\\
		&\leq |y_q|+|x_q|+4 = 2|x_q|+4. \label{eq3}
	\end{align}
	
	If $|x_q|  \leq \f-2 $, then  $d(y,\lambda^{[x;ijk]}) \leq r$, a contradiction to our assumption. So   $|x_q|=\f-1$ (because $|x_q|< \f$). 
	
	Suppose $r$ is odd. Since $|x_l|=1$,  $\sum_{s \in \{i, j, k\}\setminus \{q\}} |x_s|=\f$. If $|x_t|\geq 3$ for any $t\in \{i, j, k\}\setminus \{q\}$, then $|x_q|+|x_t|>\cil$, which is a contradiction. Therefore,  $\sum_{s \in \{i, j, k\}\setminus \{q\}} |x_s|\leq 4$. This implies that $\f\leq 4$ and thus, $r\leq 9$, a contradiction. 
	
	Suppose $r$ is even. Since $r\geq 10$, there exists a $t\in\{i,j,k\}\setminus \{q\}$ such that $|x_t|\geq 2$. This implies that $|x_q|+|x_t|>\frac{r}{2}$. This is a contradiction.
		
	Thus we conclude that $d(y, \lambda^{[x;ijk]}) \leq r$, and therefore $y \in \n[\lambda^{[x;ijk]}, \D_2]$. 
	Hence, $\n[x,\D_1]\subseteq \n[\lambda^{[x;ijk]},\D_1]$. This proves Claim \ref{claim:fourvertex2}.  
	
	
	From Proposition \ref{flag}, $\D_2\simeq \D_2 \setminus x$. Now, applying the same arguments for each $c\in C$, we find that $\D_2\simeq \mathrm{Ind}_{\D_2}(V(\D_2)\setminus C)$. Moreover, since $A=B\cup C$ and $\D\simeq \D_1\simeq \mathrm{Ind}_{\D_1}(V(\D_1)\setminus B)= \D_2\simeq \mathrm{Ind}_{\D_2}(V(\D_2)\setminus C)$, it follows that $\D\simeq \mathrm{Ind}_{\D}(V(\D)\setminus A)$. This completes the proof.
\end{proof}

\end{document}
